% Auxiliary structural lemma; not a solution of the general conjecture.
\section{First differences of products of intervals}
\begin{lemma}[Shrinking interval masses]
Let $(d_j)_{j\in\mathbb Z}$ be nonnegative and log-concave, with
interval support. For fixed integers $A,B$, put
\[
 g_n=\sum_{j=A+n}^{B-n}d_j,
\]
where an empty sum is zero. On any range where the intervals are
nonempty, $(g_n)$ is log-concave. The statement also holds after
extending it by zero past the first empty interval.
\end{lemma}
\begin{proof}
At one center write $a=A+n$, $b=B-n$, $G=\sum_a^b d_j$,
$I=\sum_{a+1}^{b-1}d_j$, $L=d_a+d_b$, and
$R=d_{a-1}+d_{b+1}$. If $a\ge b$, the next interval is empty
and the desired inequality is immediate. Otherwise $G=I+L$.
For every interior index, log-concavity and interval support give
\[
 d_{a-1}d_j\le d_a d_{j-1},\qquad
 d_{b+1}d_j\le d_b d_{j+1}.
\]
After summation, $RI\le LG$. Therefore
$G^2-(G+R)I=GL-RI\ge0$, as required. Zero entries outside
the support require no division and cause no exception.
\end{proof}

\begin{theorem}[Log-concave rising differences of interval products]
Let $C(z)=\prod_{i=1}^s[t_i]=\sum_{j=0}^D c_jz^j$, with all
$t_i\ge1$. The sequence
\[
 d_j=c_j-c_{j-1}\qquad
 (0\le j\le\lfloor D/2\rfloor),\quad c_{-1}=0,
\]
is nonnegative and log-concave, with interval positive support.
\end{theorem}
\begin{proof}
Induct on the number of factors, beginning with $C=1$.
Extend the old sequence $d$ by zero outside
$[0,\lfloor D/2\rfloor]$. For $C'=[t]C$ and
$0\le n\le\lfloor(D+t-1)/2\rfloor$, symmetry and telescoping give
\[
 d'_n=c_n-c_{n-t}
     =\sum_{j=\max(0,n-t+1)}^{\min(n,D-n)}d_j.
\]
The expression is nonnegative. For $n\le\lfloor D/2\rfloor$
it is the ordinary length-$t$ window of $d$, which is log-concave
by convolution. At the junction, the next shrinking-interval mass
is at most the ordinary next window mass; thus the same
log-concavity inequality remains valid. Explicitly, if $D=2m$,
the next mass omits $d_m$ from that window, and if $D=2m+1$,
the two next masses agree.

For subsequent centers, the summation interval has endpoints
$n-t+1$ and $D-n$, so the shrinking-interval lemma applies.
At the first such center when $D$ is odd, the preceding enlarged
interval differs from the preceding actual window only by the
zero term $d_{\lfloor D/2\rfloor+1}$. Clipping the lower endpoint
at zero is identical to zero extension. Empty final intervals
give a terminal string of zeros. The supports of the moving and
shrinking windows overlap at their junction, so the positive
support is an interval. This completes the induction.
\end{proof}

\paragraph{Limitation for the current problem.}
The analogous assertion fails for the reverse of a product of
two monotone-composition rank polynomials. For example, for
$X=(2,13,16,38)$ and $Y=(2)$, its positive first differences at
indices $29,30,31$ are $13,7,4$, and $7^2<13\cdot4$.
Consequently the theorem can be applied to individual interval
products, but not to their arm expansions without a further
argument controlling the sum.
